% =====================================================================
% Preâmbulo: formatação do relatório (base ABNT, mesmo modelo do SAACI)
% Não mexer aqui sem combinar com o grupo.
% =====================================================================

% --- Idioma e codificação ---
\usepackage[T1]{fontenc}
\usepackage[utf8]{inputenc}
\usepackage[brazilian]{babel}
\usepackage{csquotes}

% --- Fonte: Times New Roman 12 ---
\usepackage{mathptmx}
\usepackage{microtype}

% --- Página A4, margens: sup. e esq. 3 cm; inf. e dir. 2 cm ---
\usepackage[a4paper,top=3cm,left=3cm,bottom=2cm,right=2cm,headsep=0.6cm,headheight=15pt]{geometry}

% --- Espaçamento 1,5 e parágrafo com recuo de 1,25 cm, justificado ---
\usepackage{setspace}
\onehalfspacing
\usepackage{indentfirst}
\setlength{\parindent}{1.25cm}
\setlength{\parskip}{0pt}

% --- Número de página no canto superior direito ---
\usepackage{fancyhdr}
\pagestyle{fancy}
\fancyhf{}
\fancyhead[R]{\thepage}
\renewcommand{\headrulewidth}{0pt}
\fancypagestyle{plain}{\fancyhf{}\fancyhead[R]{\thepage}\renewcommand{\headrulewidth}{0pt}}

% --- Títulos de seção: "1. Introdução", em negrito ---
\usepackage{titlesec}
\renewcommand{\thesection}{\arabic{section}}
\renewcommand{\thesubsection}{\arabic{section}.\arabic{subsection}}
\renewcommand{\thesubsubsection}{\arabic{section}.\arabic{subsection}.\arabic{subsubsection}}
\titleformat{\section}{\normalfont\fontsize{13}{15}\bfseries}{\thesection.}{0.4em}{}
\titleformat{\subsection}{\normalfont\normalsize\bfseries}{\thesubsection}{0.5em}{}
\titleformat{\subsubsection}{\normalfont\normalsize\bfseries}{\thesubsubsection}{0.5em}{}
\titlespacing*{\section}{0pt}{18pt}{6pt}
\titlespacing*{\subsection}{0pt}{12pt}{4pt}
\titlespacing*{\subsubsection}{0pt}{10pt}{4pt}

% --- Figuras, tabelas e quadros: legenda ACIMA, fonte ABAIXO, tamanho 10 ---
\usepackage{graphicx}
\usepackage{float}
\usepackage{booktabs}
\usepackage{tabularx}
\usepackage[font={footnotesize,singlespacing},labelsep=endash,justification=centering,skip=4pt]{caption}
\usepackage{subcaption}
\captionsetup[figure]{name=Figura,position=above}
\captionsetup[table]{name=Tabela,position=above}
\captionsetup[sub]{labelformat=parens,labelsep=space,position=below}

% Quadro (texto, bordas fechadas): ambiente próprio, numeração própria
\usepackage{amsmath}
\usepackage{newfloat}
\DeclareFloatingEnvironment[fileext=loq,name=Quadro,placement=H]{quadro}
\captionsetup[quadro]{position=above}

% Diagramas desenhados no próprio LaTeX (câmera, volume de visão, Phong, retas, scan line, baricêntrica)
\usepackage{tikz}
\usetikzlibrary{arrows.meta,calc}

% \fonte{...}: linha "Fonte: ..." abaixo de figura, tabela ou quadro
\newcommand{\fonte}[1]{\par\vspace{2pt}{\centering\footnotesize Fonte: #1\par}}

% \figura{arquivo}{largura}{legenda}{rótulo}{fonte}
\newcommand{\figura}[5]{%
  \begin{figure}[H]
    \centering
    \caption{#3}\label{#4}
    \includegraphics[width=#2]{#1}
    \fonte{#5}
  \end{figure}}

% Caminho das imagens
\graphicspath{{figuras/}}

% --- Listas ---
\usepackage{enumitem}
\setlist{nosep,leftmargin=1.25cm}

% --- Citação direta longa: recuo de 4 cm, fonte 10, espaço simples, sem aspas ---
\newenvironment{citacao}{%
  \par\vspace{6pt}\begin{list}{}{\setlength{\leftmargin}{4cm}\setlength{\rightmargin}{0pt}}
  \item[]\footnotesize\singlespacing}{%
  \end{list}\vspace{6pt}}

% --- Referências ABNT (NBR 6023:2018) com biblatex-abnt ---
\usepackage[style=abnt,backend=biber,justify=false,extrayear]{biblatex}
\addbibresource{referencias.bib}
% NBR 10520:2023: no texto, autor só com a inicial maiúscula -> (Martin, 2017).
% Na lista de referências o SOBRENOME continua em CAIXA ALTA.
\makeatletter
\renewcommand*{\UpperOrSCCite}[1]{#1}
\DeclareFieldFormat{uppercasecite}{#1}
\makeatother
\setlength\bibitemsep{\baselineskip}
\renewcommand*{\bibfont}{\normalsize\singlespacing}

% --- Links (sem caixas coloridas) ---
\usepackage[hidelinks]{hyperref}
\urlstyle{same}

% --- Comando para marcar o que falta confirmar (aparece em vermelho) ---
\usepackage{xcolor}
\newcommand{\confirmar}[1]{\textcolor{red}{[CONFIRMAR: #1]}}
% Para a versão final: descomente a linha abaixo para esconder as marcações
% \renewcommand{\confirmar}[1]{}
